\documentclass[11pt]{article}
\usepackage[utf8]{inputenc}
\usepackage[T1]{fontenc}
\usepackage{lmodern}
\usepackage[margin=1in]{geometry}
\usepackage{csquotes}
\usepackage[style=authoryear,backend=biber,sorting=nyt,maxbibnames=99,giveninits=true,uniquename=init]{biblatex}
\usepackage[hidelinks]{hyperref}
\addbibresource{references.bib}
\begin{document}
\section*{Long-Term Monitoring and Research Needs}

\subsection{Overview}

A recurring theme across the subsurface (tile) drainage literature is that the current evidence base remains provisional: many headline results are drawn from short records, single sites, or specific agro-ecological conditions, and their authors explicitly caution against premature generalisation. The dominant frontier message is therefore not a settled conclusion but a research agenda---one that calls for extended monitoring, broader spatial and geographic replication, harmonised and standardised methods, field-scale validation of models and emerging technologies, and stronger links between environmental, agronomic, economic, and social evidence. Several authors frame these needs at the level of the whole discipline: \textcite{Rogger2017} argue for adopting complex-systems thinking to link processes across time scales, conducting long-term experiments on physical--chemical--biological process interactions, focusing on connectivity and patterns across spatial scales, developing typologies and meta-analyses to generalise findings, strengthening multidisciplinary research consortia with shared data and terminology, and using long-term observations under controlled or known boundary conditions wherever possible. This section synthesises the specific monitoring and research needs identified across the reviewed studies, organised by the temporal, spatial, methodological, integrative, and programmatic dimensions along which the field is asked to advance.

The scale of the outstanding uncertainty is captured by the frequency of direct appeals for more work. Studies conclude, variously, that ``more research is needed'' and that ``more verification would be beneficial'' \parencite{SHALIT1995}, that ``the complexity of this subject requires more research'' \parencite{Malone2014}, that ``further research in this area is needed'' \parencite{Woo2023}, and that ``further evaluations are warranted before any concrete conclusions could be drawn'' \parencite{GUPTA1993}. These are not rhetorical closings but signposts to genuine gaps in evidence, method, and synthesis that the subsections below make explicit.

\subsection{The Case for Long-Term, Multi-Year Monitoring}

The single most widely voiced need is for records that are long enough to separate genuine treatment effects from the natural variability of weather, hydrology, and crop growth. \textcite{Clifford2023} emphasise the need for truly long-term experiments and monitoring to accurately assess management impacts, and \textcite{BUHLER1993} caution that a lag time of several years may be necessary before the impact of altered practices can be objectively evaluated. \textcite{Ali2013} judge that longer monitoring periods---at least five years---and continued data collection are desirable, and warn that long-term consequences may need continued observation if drainage systems are expanded in length or depth over time. In Egypt, \textcite{Ali2001} report that monitoring work will continue, that within a few years more and longer time series will become available for the whole country, and that a more detailed analysis is expected to become possible as those records lengthen.

At the site scale, numerous studies attach their conclusions to explicit requirements for additional years of observation. \textcite{Jaynes2008} state that continued monitoring of the denitrification-wall treatment for additional years is necessary to clearly quantify its effectiveness given natural variability in weather and crop growth, that extended monitoring is required to quantify effectiveness over time, and that further work is needed to determine the long-term cost per unit of nitrate removed and the system's potential as a management tool. \textcite{Dalhem2019} note that continued sampling and measurements will reveal how long the positive effects of the treatments persist. \textcite{Osterholm2005} judge it worthwhile to monitor a test site for several more years to better assess the functionality of a combination of two techniques, while \textcite{ROLLIN1987} plan to gather additional observations covering a longer service-life period. Similar calls come from \textcite{Wesström2015}, who require verification over a longer monitoring period, \textcite{Reinhart2020}, who call for more widespread and longer-duration evaluations to better understand a practice's potential impact, and \textcite{Rijal2012}, who identify a need to assess the replicability of results in time. \textcite{GiménezValero2025} similarly conclude that additional multi-year measurements will be necessary to confirm whether observed differences persist under comparable environmental and management conditions. Underlying all of these is the disciplinary-scale argument of \textcite{SPALING1995} for long-term monitoring of environmental change at various scales.

\subsection{Temporal Resolution and Seasonal Coverage}

Beyond record length, several authors highlight gaps in the temporal \emph{resolution} and seasonal coverage of monitoring. \textcite{Klaiber2020} make this explicit, calling for additional years of monitoring, improved monitoring during the non-growing season, and year-round monitoring to better understand seasonal phosphorus exports and develop more robust conclusions. The importance of avoiding gaps in high-frequency records is stressed by \textcite{Hodaj2017}, who advise that future studies check instruments as frequently as possible to limit gaps in the data collected. \textcite{Wesström2015} add a timing consideration specific to infrastructure interventions, noting that it is more appropriate to monitor leaching of nutrients after drain repair.

\subsection{Spatial Replication and Geographic Diversity}

A second major frontier is the need to test whether site-specific findings hold across other locations, soils, climates, and management regimes. \textcite{Menegaz2025} state that additional studies are needed before broad adoption of their results can be encouraged beyond the local conditions of their study, and \textcite{Zimmerman2017b} call for similar data to be reported for other sites and years. \textcite{Ali2001} plan to collect more data from other sample areas in Egypt reflecting the different agro-ecological conditions prevalent in the country. \textcite{Hodaj2017} argue that further exploration of the two-stage ditch in other settings would benefit research on the practice, and \textcite{Bigelow2001} recommend that further research address effects where putting greens are severely compacted, possess low percolation rates, or are maintained in unfavourable environments such as shade or excessive wetness, while suggesting the technology may have merit under other situations. \textcite{Terrado2015} likewise conclude that additional monitoring and research supporting their ideas and findings is warranted. At the level of evidence synthesis, \textcite{Wu2025} argue that future syntheses will require more robust studies that are geographically diverse and methodologically harmonized, and that with more extensive experimental evidence future meta-analyses will be better positioned to disentangle management effects.

\subsection{Field-Scale Validation of Models and Emerging Technologies}

Much of the reviewed work rests on models, pilot studies, or laboratory-scale results whose authors explicitly defer confident application until field-scale validation is available. \textcite{Mitchell2025} state plainly that ``field-scale research is needed to validate these results'' and that ``further research to validate our models is needed before they are widely applied.'' \textcite{Kalcic2018} similarly conclude that continued efforts to validate practice effectiveness at the field and reach scales will be needed to verify if models are correct, and that these areas are ripe for future modeling development and validation. \textcite{Garmdareh2018} recommend more rigorous testing of a model before its extensive use, and \textcite{Yu2016} call for validating their conclusions with more experiments. \textcite{SHALIT1995} identify a specific validation target, arguing that additional investigations should determine how well the leaching-fraction model applies to predicting concentrations from preferential-flow situations. For emerging technologies, \textcite{DarziNaftchali2018b} stress that before new technologies can be implemented on a large scale, their environmental effects need to be quantified at the level of pilot areas, and \textcite{GUPTA1993} require further measurements and evaluations before concrete conclusions can be drawn.

\subsection{Monitoring Infrastructure, Instrumentation, and Data Tools}

A distinct set of needs concerns the practical apparatus of monitoring---the samplers, loggers, diagnostic methods, and software that make long-term data collection feasible and reliable. \textcite{Kanwar1999} describe how automatic samplers can be connected to a monitoring system to collect discrete water samples alongside continuous samples, how additional equipment can be connected to data loggers to measure soil temperature or soil moisture, and how relays or pressure switches could replace electronic transmitters to detect sump-pump operation; they further advise minimising the distance between monitoring equipment and data loggers, downloading data automatically at set intervals to minimise data loss, and performing annual maintenance including cleaning check valves and calibrating flowmeters. Diagnostic methods themselves are a research target: \textcite{Petersen2012} propose reducing pressure variations by doubling the equipment and injecting smoke from both ends of a drain section, potentially enabling work with longer drain lines, while noting that further studies are needed to discover the details of the pathways and to disclose reasons for observed clustering. \textcite{MartínezBeltrán2007} recommend observing selected pipes to determine the state of the envelope and perforations, observing the surrounding soil because of its influence on radial flow in the drain trench, and permanently monitoring the sampling process to verify that favourable results are maintained over time.

Software and data infrastructure form an increasingly prominent part of this frontier. \textcite{Shedekar2018} advise readers to explore data products and tools further, to visit the National Elevation Dataset and USGS websites for the latest data availability and updates, and to track ongoing developments of these tools, while also arguing that improved drainage-management practices are needed to increase adaptation to a changing climate. A series of tool-development efforts describe themselves as explicitly ongoing and open to community input: \textcite{Falasy2025} note that development of their tools is continuing---spanning creation of topologically sound, one-way-flow tile networks and application tools for configuring the burying and sizing of drainage networks---and welcome suggestions and testing with datasets beyond those provided; \textcite{Falasy2025b} report an ongoing development process with numerous revisions based on test-user feedback and encourage application to data beyond the samples provided; and \textcite{Falasy2024} state that development of their tools will be continual and encourage testing with other data.

\subsection{Standardisation, Harmonisation, and Comparability of Data}

For the accumulating body of monitoring data to support robust synthesis, several authors identify a need for standardised and harmonised methods. \textcite{Wu2025} tie the value of future meta-analyses to methodologically harmonized studies and comprehensive, long-term monitoring of the full suite of environmental impacts. \textcite{Zhang2018} set out a detailed programme for environmental indicators: improving the quality and comparability of existing indicators over time through harmonized methods for data collection, monitoring, interpretation, and reporting; developing new indicators that address emerging environmental and policy concerns while ensuring scientific credibility, sensitivity to change, and ease of interpretation; developing standardized procedures or guidelines for evaluation, including standards, thresholds, reference points, and complementary information for validation; recognising data gaps and improving methodologies for data upscaling, uncertainty assessment, and analysis of contradictory indicators; and improving or adapting environmental models for indicator modeling while incorporating earth-observation data and reducing barriers to remote-sensing data access. These priorities complement the call by \textcite{Rogger2017} for typologies, meta-analyses, and shared data and terminology across research consortia.

\subsection{Understanding Complex Processes, Pathways, and Interactions}

A further frontier concerns mechanistic understanding: identifying the pathways, relationships, and process interactions that current data only partially resolve. \textcite{Malone2014} conclude simply that the complexity of the subject requires more research, and \textcite{Cui2016} state that more research is needed to define the relationships they observed. \textcite{Petersen2012} identify unresolved transport pathways and clustering patterns requiring further investigation. \textcite{Tomer2008} intend to pursue their central idea as an avenue of research, to undertake bacterial source tracking in future work within their watershed, and note that a specific transport pathway remains a subject for future research. \textcite{Ali2013} call for a separate geochemical modelling study to validate whether geochemical reactions bind free heavy metals with the solid matrix. Process-level questions extend to plant response: \textcite{GiménezValero2025} note that an observation consistent with the role of phosphorus in flowering merits further study under controlled conditions, and \textcite{Tamura2021} report that further experiments are required to clarify a specific point in more detail. \textcite{LAL1993} indicate that results of experiments on plant-nutrient relations will be presented in a companion report and that a long-term experiment is being analysed to test their hypotheses. These needs echo the overarching argument of \textcite{Rogger2017} for long-term experiments on physical--chemical--biological process interactions and for a focus on connectivity and patterns across spatial scales.

\subsection{Integrating Environmental, Agronomic, Economic, and Social Evidence}

Several authors argue that monitoring must move beyond environmental endpoints alone to encompass agronomic, economic, and social dimensions, and to connect with decision-making. \textcite{Zhang2018} call for improved linkages of environmental indicators with indicators tracking economic and social change, in terms of both causes and effects. \textcite{Zhu2025} require further assessment to evaluate actual growth in conjunction with economic, social, and environmental factors, and \textcite{Brikmans2025} argue that future investigations should elucidate strategies for prolonging physical benefits while critically assessing the agronomic and economic feasibility of such approaches. \textcite{Monaghan2016} identify a need for further research on the practical implications of management changes---in their case, having cows off-pasture for extended periods. At the science--policy interface, \textcite{Letey1999} argue that plans should be updated as monitoring, additional studies, and local actions reveal new facts, and that the scientific community needs to evaluate its role in policymaking and focus research on questions of greatest value to decision makers as well as to scientific peers. \textcite{Clifford2023} call for scientific assessment and prioritisation of which approaches have the largest effects under what circumstances, alongside truly long-term experiments and monitoring. \textcite{Kalcic2018} likewise argue that further work is needed to assess current knowledge and prioritise knowledge gaps, that clear messaging through education and extension will be imperative to achieve implementation at the scales needed to treat water-quality needs in the Mississippi River Basin and Great Lakes regions, and that these areas are ripe for future modeling development and validation.

\subsection{Ongoing and Planned Monitoring Programmes}

Finally, the literature documents a number of monitoring programmes that are explicitly ongoing or planned, underscoring that many conclusions are interim by design. \textcite{Tournebize2012} report that in 2011 a performance-assessment programme was established to determine whether co-constructed wetlands comply with new regulations, that intensive hydrological monitoring will continue to assess the implemented system's performance, and that resizing will be considered if needed in accordance with current regulations. \textcite{Kjaersgaard2016} note that monitoring was to begin in spring 2016, and \textcite{Ali2001} confirm that monitoring work will continue. \textcite{Ohlendorf2002} describe how data from a monitoring programme were used to identify needs for contingency plans, to recommend further research and monitoring to improve site management, and how further studies were subsequently conducted to resolve unknowns identified by a risk assessment. Other authors flag forthcoming results: \textcite{Drury2014} indicate that data from remaining treatments will be presented in a subsequent paper. Additional calls for continued or expanded empirical effort come from \textcite{Rekha2011}, who state that more treatment studies are required, and from \textcite{Ogino2007}, who express hope that research and development in the field will be actively promoted. Collectively, these programmes and appeals reinforce the central frontier message: robust, generalisable conclusions about subsurface drainage await the longer, broader, better-harmonised, and more integrative evidence base that the field is actively working to build.

\printbibliography
\end{document}
